% !TEX root = ../main.tex

\begin{figure}[H]
  \centering
  \resizebox{\textwidth}{!}{%
    \begin{tikzpicture}[
      % Sin `text width`: el nodo se ajusta a su rotulo y asi ninguno se parte
      % con guion. Con 3,7 cm entre lineas de vida caben todos holgados.
      papel/.style={nodoDiag, minimum height=0.9cm, font=\small},
      papelFin/.style={papel, draw=diagFocoBorde, fill=diagFocoFondo, very thick},
      lineaVida/.style={densely dotted, thick, draw=diagNodoBorde},
      mensaje/.style={-{Latex[length=2mm]}, semithick},
      devuelve/.style={-{Latex[length=2mm]}, semithick, densely dashed},
      mensajeFin/.style={-{Latex[length=2mm]}, very thick, draw=diagFocoBorde},
      % Relleno opaco: sin el, las lineas de vida punteadas atraviesan los
      % rotulos largos y los vuelven ilegibles.
      etiqueta/.style={font=\small, midway, above, inner sep=1.5pt, fill=white},
      propia/.style={font=\small, anchor=west, inner sep=1.5pt, fill=white},
      bloque/.style={draw=diagGrupoBorde, semithick, densely dashed,
        fill=diagGrupo, fill opacity=0.4, text opacity=1,
        rounded corners=2pt},
      rotuloBloque/.style={font=\small\bfseries, anchor=north west,
        inner sep=2.5pt, text=diagGrupoBorde!65!black, fill=white},
      apunte/.style={font=\small, align=center, draw=diagNodoBorde,
        fill=white, inner sep=3pt, rounded corners=2pt},
      ]

      % --- Participantes y lineas de vida ---
      % Centrados en -1,3 para que su borde inferior quede justo encima del
      % arranque de las lineas de vida, sin tener que mover ningun mensaje.
      \node[papel]    at (0,-1.3)    {El fragmentador};
      \node[papel]    at (3.7,-1.3)  {La colección};
      \node[papel]    at (7.4,-1.3)  {El agrupador};
      \node[papel]    at (11.1,-1.3) {El troceador};
      \node[papelFin] at (14.8,-1.3) {Lo que se indexa};
      \foreach \x in {0, 3.7, 7.4, 11.1, 14.8}{
        \draw[lineaVida] (\x,-1.75) -- (\x,-12.6);
      }

      % --- Lectura de la coleccion ---
      \draw[mensaje]  (0,-2.4) -- node[etiqueta]{lee 837 registros} (3.7,-2.4);
      \draw[devuelve] (3.7,-3.0) -- node[etiqueta]{la colección} (0,-3.0);

      % --- Deduplicacion ---
      \draw[mensaje] (0,-3.9) -- node[etiqueta]{las 288 guías} (7.4,-3.9);
      \draw[mensaje] (7.4,-4.5) -- ++(0.8,0) -- ++(0,-0.45) -- ++(-0.8,0);
      \node[propia] at (8.3,-4.75) {agrupa por (nombre, contenido)};
      % «de guia» no sobra: el corpus entero tiene 398 unidades, de las que 210
      % son de guia. Sin el matiz, esta cifra parece contradecir la otra.
      \draw[devuelve] (7.4,-5.6)
      -- node[etiqueta]{210 unidades de guía: 78 copias que no se repiten} (0,-5.6);

      % --- Listados compuestos ---
      \draw[mensaje] (0,-6.2) -- ++(0.8,0) -- ++(0,-0.45) -- ++(-0.8,0);
      \node[propia] at (0.9,-6.45)
      {compone los listados: plan, catálogo, ficha y mención};

      % --- El bucle, una unidad cada vez ---
      \node[bloque, fit={(-1.4,-7.3) (16.2,-10.5)}] (b1) {};
      \node[rotuloBloque] at (b1.north west) {por cada unidad};
      \node[font=\small, anchor=north west, fill=white] at (3.4,-7.42)
      {[210 de guía, 8 de salidas y los listados compuestos]};
      \draw[mensaje] (0,-8.5) -- node[etiqueta]{encabezado + texto} (11.1,-8.5);
      \draw[mensaje] (11.1,-9.1) -- ++(0.8,0) -- ++(0,-0.45) -- ++(-0.8,0);
      \node[propia] at (12.0,-9.35) {corta, empaqueta y fusiona lo corto};
      \draw[devuelve] (11.1,-10.1)
      -- node[etiqueta]{fragmentos numerados} (0,-10.1);

      % --- Lo que se escribe, y lo que se escribe sin trocear ---
      \draw[mensajeFin] (0,-11.3) -- node[etiqueta]{1.413 fragmentos} (14.8,-11.3);
      \draw[mensaje]    (0,-12.0)
      -- node[etiqueta]{86 avisos de «sin guía»} (14.8,-12.0);
      \node[apunte] at (7.4,-13.1)
      {los 86 avisos son el único mensaje que no pasa por el troceador:\\
        no hay texto que cortar, pero la asignatura no desaparece del corpus};
    \end{tikzpicture}%
  }
  \caption[Secuencia del proceso de fragmentación]{Diagrama de secuencia del proceso de fragmentación. Ilustra la deduplicación inicial de guías compartidas, la composición de listados, el particionado recursivo por unidad documental y la emisión directa de los fragmentos informativos para materias sin guía docente. Fuente: Elaboración propia.}

  \label{fig:secuencia_chunking}
\end{figure}
